\chapter{Point Estimation: Unbiasedness and Efficiency}
\label{ch:estimation}

The previous chapter introduced point estimators and the notions of bias and
variance. This chapter studies them in depth. We first build estimators from data
by matching sample quantities to model quantities, then ask which estimators are
\emph{unbiased}, see how nonlinearity creates bias, and finally compare
unbiased estimators by their variance to find the most \emph{efficient}.

\section{Estimators Built from Data}

A natural way to estimate a parameter is to match a population quantity—a mean, a
probability—to its sample counterpart. The law of large numbers is what makes
this work.

\begin{example}[A binomial parameter from queue data]
At a station, the number of women in each of $100$ queues of length $10$ is
recorded. Modeling the count in a queue as $\Bin(10,p)$, estimate $p$.

\solution If each of the $10$ positions is independently a woman with probability
$p$, the count is $\Bin(10,p)$ with mean $10p$. The sample mean of the counts is
$\xbar=4.35$, which we equate to $10p$, giving $\hat p=0.435$. We matched the
model's mean to the data's mean.
\end{example}

\begin{example}[Two estimators for a geometric parameter]
The number of cycles to pregnancy for $93$ women is recorded, modeled as
$\Geo(p)$. Propose estimators of $p$.

\solution Two natural choices arise. Since $\E[X]=1/p$, we can estimate $\E[X]$
by the sample mean and invert: $\hat p=1/\xbar=93/331$. Alternatively, since
$p=\Prob(X=1)$, we can use the observed fraction of women who conceived in the
first cycle, $\hat p=29/93$. The two estimators—$0.281$ and $0.312$—need not
agree; choosing between them is exactly the kind of question this chapter
addresses. (For the $474$ nonsmoking women, the analogous estimates are
$474/1285=0.369$ and $198/474=0.418$.)
\end{example}

\begin{example}[Was London bombed at random? A Poisson fit]
An area of South London was divided into $576$ quarter-kilometer squares, and the
number of V-bomb hits in each was tallied:
\[
\begin{array}{c|cccccccc}
\text{hits} & 0 & 1 & 2 & 3 & 4 & 5 & 6 & 7\\\hline
\text{squares} & 229 & 211 & 93 & 35 & 7 & 0 & 0 & 1
\end{array}
\]
If hits fall purely at random, the count per square should be Poisson. Estimate
its mean and assess the fit.

\solution Modeling the count as $\Pois(\mu)$ with $\E[X]=\mu$, we estimate $\mu$
by the average number of hits per square,
\[
  \hat\mu=\frac{0\cdot 229+1\cdot 211+2\cdot 93+\cdots+7\cdot 1}{576}\approx 0.93.
\]
Plugging $\hat\mu$ into the Poisson pmf and comparing predicted to observed
relative frequencies:
\[
\begin{array}{c|ccccc}
\text{hits} & 0 & 1 & 2 & 3 & 4\\\hline
\text{observed} & 0.40 & 0.37 & 0.16 & 0.06 & 0.01\\
\text{Poisson} & 0.39 & 0.37 & 0.17 & 0.05 & 0.01
\end{array}
\]
The agreement is remarkable—evidence that the bombing was indeed random rather
than precisely targeted. A formal version of this comparison, the chi-square
goodness-of-fit test, appears in Chapter~\ref{ch:gof}.
\end{example}

\section{Unbiased Estimators}

Recall that $\hat\theta_n$ is \emph{unbiased} for $\theta$ if
$\E[\hat\theta_n]=\theta$. Unbiasedness means the estimator is correct
\emph{on average}, with no systematic over- or under-shooting.

\begin{example}[An unbiased estimator of $\theta^2$]
Let $X_1,\dots,X_n$ be a sample from the uniform distribution on $[-\theta,\theta]$.
Show that $T=\tfrac3n\sum_{i=1}^n X_i^2$ is unbiased for $\theta^2$.

\solution For a $\Unif(-\theta,\theta)$ variable, $\E[X_i]=0$ and
$\Var(X_i)=\tfrac{(2\theta)^2}{12}=\tfrac{\theta^2}{3}$, so
$\E[X_i^2]=\Var(X_i)+\E[X_i]^2=\tfrac{\theta^2}{3}$. By linearity,
\[
  \E[T]=\frac3n\sum_{i=1}^n\E[X_i^2]=\frac3n\cdot n\cdot\frac{\theta^2}{3}=\theta^2. \qedhere
\]
\end{example}

\begin{example}[An unbiased estimator of a probability]
For a $\Geo(p)$ sample, estimate $\Prob(X\le 3)$ by the relative frequency
$S=\tfrac1n\#\{i:X_i\le 3\}$. Show $S$ is unbiased.

\solution Introduce indicators $Y_i=\mathbf 1\{X_i\le 3\}$, so
$\E[Y_i]=\Prob(X_i\le 3)$. Then
\[
  \E[S]=\frac1n\sum_{i=1}^n\E[Y_i]=\Prob(X\le 3).
\]
A relative frequency is always an unbiased estimator of the corresponding
probability—no matter the underlying distribution.
\end{example}

\begin{remark}[Why the sample variance divides by $n-1$]
We can now explain the divisor $n-1$ promised in Chapter~\ref{ch:eda}. For an
i.i.d.\ sample with variance $\sigma^2$, the estimator
$\frac1{n-1}\sum(X_i-\Xbar_n)^2$ is exactly unbiased for $\sigma^2$, whereas
dividing by $n$ would systematically underestimate it (because the deviations are
taken about the sample mean, which is itself pulled toward the data). The
correction factor $\tfrac{n}{n-1}$ removes this downward bias.
\end{remark}

\section{Bias from Nonlinearity}

Unbiasedness is fragile under nonlinear transformation: even if $T$ is unbiased
for $\theta$, $g(T)$ is usually \emph{biased} for $g(\theta)$. Jensen's
inequality (Chapter~\ref{ch:limits}) tells us the \emph{direction} of the bias.

\begin{example}[A square root introduces bias]
With $T=\tfrac3n\sum X_i^2$ unbiased for $\theta^2$ as above, is $\sqrt T$
unbiased for $\theta$?

\solution The function $g(x)=\sqrt x$ is concave, so by Jensen's inequality
$\E[\sqrt T]\le\sqrt{\E[T]}=\sqrt{\theta^2}=\theta$. Thus $\sqrt T$ has
\emph{negative} bias: it tends to underestimate $\theta$.
\end{example}

\begin{example}[Inverting the sample mean]
For a $\Geo(p)$ sample, the law of large numbers suggests $T=1/\Xbar_n$ as an
estimator of $p$ (since $\E[X]=1/p$). Is it unbiased?

\solution The function $g(x)=1/x$ is convex on $(0,\infty)$, so Jensen's
inequality gives $\E[1/\Xbar_n]>1/\E[\Xbar_n]=p$. Hence $T$ has \emph{positive}
bias—it tends to overestimate $p$. Inverting an unbiased estimator does not
preserve unbiasedness.
\end{example}

\begin{example}[An exact bias calculation]
For a $\Pois(\mu)$ sample, the probability of zero arrivals is $e^{-\mu}$, which
one might estimate by $T=e^{-\Xbar_n}$. Compute $\E[T]$.

\solution Write $T=e^{-Z/n}$ with $Z=\sum X_i\sim\Pois(n\mu)$. Then
\[
  \E[T]=\sum_{k=0}^\infty e^{-k/n}\,\frac{(n\mu)^k}{k!}e^{-n\mu}
   =e^{-n\mu}\sum_{k=0}^\infty\frac{(n\mu e^{-1/n})^k}{k!}
   =e^{-n\mu}\,e^{\,n\mu e^{-1/n}}=e^{-n\mu(1-e^{-1/n})}.
\]
Since $1-e^{-1/n}\neq 1/n$, this is not equal to $e^{-\mu}$: the estimator is
biased, though the bias vanishes as $n\to\infty$.
\end{example}

\section{Efficiency: Comparing Unbiased Estimators}

Among unbiased estimators, the better one has the smaller variance (equivalently,
since bias is zero, the smaller MSE). Such an estimator is more \emph{efficient}.

\begin{example}[Optimally combining two sample means]
Two independent samples, of sizes $n$ and $m$ from the same distribution with
mean $\mu$, give sample means $\Xbar_n$ and $\Ybar_m$. Combine them as
$T=r\Xbar_n+(1-r)\Ybar_m$. (a)~Show $T$ is unbiased. (b)~Find the $r$ minimizing
its variance.

\solution (a)~By linearity, $\E[T]=r\mu+(1-r)\mu=\mu$ for every $r$, so $T$ is
unbiased. (b)~By independence,
\[
  \Var(T)=r^2\frac{\sigma^2}{n}+(1-r)^2\frac{\sigma^2}{m}.
\]
Differentiating in $r$ and setting the derivative to zero gives
$\tfrac{2r}{n}-\tfrac{2(1-r)}{m}=0$, i.e.\ $r=\tfrac{n}{n+m}$. The optimal weights
are proportional to the sample sizes—the larger sample is trusted more.
\end{example}

\begin{example}[Pooling two variance estimators]
Independent normal samples of sizes $n$ and $m$ have sample variances $S_X^2$
and $S_Y^2$, both unbiased for the common $\sigma^2$, with
$\Var(S_X^2)=\tfrac{2\sigma^4}{n-1}$ and $\Var(S_Y^2)=\tfrac{2\sigma^4}{m-1}$.
Among estimators $aS_X^2+bS_Y^2$, find the best.

\solution For unbiasedness, $\E[aS_X^2+bS_Y^2]=(a+b)\sigma^2=\sigma^2$ forces
$a+b=1$. Writing $b=1-a$,
\[
  \Var(aS_X^2+(1-a)S_Y^2)=2\sigma^4\Bigl(\frac{a^2}{n-1}+\frac{(1-a)^2}{m-1}\Bigr).
\]
Minimizing over $a$ gives $a=\tfrac{n-1}{n+m-2}$ and hence
$b=\tfrac{m-1}{n+m-2}$. This is the \emph{pooled} variance estimator, weighting
each sample by its degrees of freedom—the standard estimator behind the
two-sample $t$-test.
\end{example}

The maximum likelihood method of the next chapter provides a systematic recipe
for constructing estimators that, under mild conditions, are asymptotically
unbiased and asymptotically the most efficient of all.
